En la figura 1 adjunta, hem situat una càrrega puntual $Q = +1,00\ \mathrm{nC}$ a l'origen de coordenades.

\figura[0.6]{fig1.png}
\begin{center}\small Figura 1\end{center}

\begin{apartat}{1,25}
Determineu a quina distància de la càrrega el potencial és igual a 3,00~V, 6,00~V, 9,00~V i 12,0~V, respectivament. Representeu dins de la figura les línies equipotencials corresponents als potencials de 3,00~V, 6,00~V, 9,00~V i 12,0~V. Digueu si la distància entre les línies equipotencials és constant i quant val o valen aquestes distàncies.
\end{apartat}

\begin{apartat}{1,25}
En la mateixa figura, representeu 8 línies de camp elèctric. Quin angle formen les línies de camp amb les línies equipotencials en el punt on es creuen?
\end{apartat}

\dades{%
  $k = \dfrac{1}{4\pi\varepsilon_0} \cong 9,00\times 10^{9}\ \mathrm{N\,m^2\,C^{-2}}$.}
